\section{Discrete Logarithm}\label{sec:discrete-log}

\subsection{The Discrete Logarithm Problem}

\begin{definition}[Group Generation Algorithm]
  A \emph{group generation\index{group generation algorithm@Group generation algorithm|textbf} algorithm} $\grgen$ is a \ppt algorithm that takes the security parameter\index{security parameter@Security parameter} $1^\secpar$ as input and outputs a group description $(\GG, p, g)$, where $\GG$ is a cyclic group of prime order $p$ with $p \geq 2^{\secpar-1}$, and $g$ is a generator of $\GG$.
\end{definition}

\begin{remark}
  The group operation is denoted multiplicatively ($g\cdot g = g^2$).
\end{remark}


\begin{figure}[tbhp]
  \begin{center}
    \begin{tcolorbox}[width=5cm]
      \begin{pchstack}[center]
        \procedure[headlinesep=1pt]{$\game{\Game~\algo{DL}}{\grgen}$}{%
          \gparam \gets \grgen(1^\secpar) \\
          x \sample \ZZ_p \\
          X \defeq g^x \\
          x' \gets \adv(\gparam, X) \\
          \pcreturn x' = x
        }
      \end{pchstack}
    \end{tcolorbox}
  \end{center}
  \caption{The discrete logarithm game\label{fig:dl-game}}
\end{figure}

\begin{definition}[Discrete Logarithm Assumption]
  Let $\grgen$ be a group generation algorithm, and let $\game{\Game~\algo{DL}}{\grgen}$ be as defined in \autoref{fig:dl-game}.
  The discrete logarithm\index{discrete logarithm dl@Discrete logarithm (DL)|textbf} (DL) problem is hard for $\grgen$ if for any \ppt algorithm $\adv$,
  \[
  \advantage{DL}{\adv, \grgen} \defeq \pr{\game{\Game~\algo{DL}}{\grgen} = 1} = \negl.
  \]
\end{definition}

\subsection{Pedersen Commitments}

\begin{definition}[Pedersen Commitment Scheme]
  Let $\grgen$ be a group generation\index{group generation algorithm@Group generation algorithm} algorithm.
  The Pedersen\index{pedersen commitment@Pedersen commitment|textbf} commitment\index{commitment scheme@Commitment scheme} scheme $\algo{PedCom}[\grgen]$ is defined as follows:
  \begin{itemize}
  \item $\algo{Setup}(1^\secpar) \rightarrow \params$: Run $\gparam \gets \grgen(1^\secpar)$, sample $h \sample \GG \setminus \{1_\GG\}$, and output $\params = (\gparam, h, \mathcal{M} \defeq \ZZ_p, \mathcal{R} \defeq \ZZ_p)$.
  \item $\algo{Commit}(\params, m, r) \rightarrow C$: Parse $\params = ((\GG, p, g), h, \mathcal{M}, \mathcal{R})$. For message $m \in \mathcal{M}$ and randomness $r \in \mathcal{R}$, output $C = g^m h^r$.
  \end{itemize}
\end{definition}

\begin{remark}
  Pedersen commitments are \emph{homomorphic\index{homomorphic commitment@Homomorphic commitment|textbf}}: for commitments $C_1 = \algo{Commit}(\params, m_1, r_1)$ and $C_2 = \algo{Commit}(\params, m_2, r_2)$, we have
  \[
  C_1 \cdot C_2 = g^{m_1 + m_2} h^{r_1 + r_2} = \algo{Commit}(\params, m_1 + m_2, r_1 + r_2).
  \]
\end{remark}

\begin{theorem}[Pedersen Commitments are Binding]\label{thm:pedersen-binding}
  Let $\grgen$ be a group generation algorithm for which the discrete logarithm\index{discrete logarithm dl@Discrete logarithm (DL)} problem is hard, then the Pedersen commitment scheme $\algo{PedCom}[\grgen]$ is binding\index{binding@Binding}.
  More precisely, for any \ppt adversary $\adv$ against the binding property of $\algo{PedCom}[\grgen]$, there exists a \ppt adversary $\bdv$ against the discrete logarithm problem such that
  \[
  \advantage{Bind}{\adv, \algo{PedCom}} = \advantage{DL}{\bdv, \grgen}.
  \]
\end{theorem}

The proof is left as an exercise (see \autoref{ex:pedersen-binding}).

\begin{theorem}[Pedersen Commitments are Hiding]\label{thm:pedersen-hiding}
  The Pedersen commitment scheme $\algo{PedCom}[\grgen]$ is perfectly hiding\index{hiding@Hiding}\index{hiding@Hiding!perfectly hiding@perfectly hiding}.
  More precisely, for any (possibly unbounded) adversary $\adv$,
  \[
  \advantage{Hid}{\adv, \algo{PedCom}} = 0.
  \]
\end{theorem}

The proof is left as an exercise (see \autoref{ex:pedersen-hiding}).

\subsection{The DDH Assumption}

\begin{figure}[tbhp]
  \begin{center}
    \begin{tcolorbox}[width=6cm]
      \begin{pchstack}[center]
        \procedure[headlinesep=1pt]{$\game{\Game~\algo{DDH}}{\grgen}$}{%
          \gparam \gets \grgen(1^\secpar) \\
          a, b \sample \ZZ_p \\
          A \defeq g^a, B \defeq g^b \\
          d \sample \{0, 1\} \\
          \pcif d = 0 \pcthen \\
          \t c \sample \ZZ_p \\
          \t C \defeq g^c \\
          \pcelse \\
          \t C \defeq g^{ab} \\
          d' \gets \adv(\gparam, A, B, C) \\
          \pcreturn d' = d
        }
      \end{pchstack}
    \end{tcolorbox}
  \end{center}
  \caption{The decisional Diffie-Hellman game\label{fig:ddh-game}}
\end{figure}

\begin{definition}[Decisional Diffie-Hellman Assumption]
  Let $\grgen$ be a group generation\index{group generation algorithm@Group generation algorithm} algorithm, and let $\game{\Game~\algo{DDH}}{\grgen}$ be as defined in \autoref{fig:ddh-game}.
  The decisional Diffie\index{decisional diffiehellman ddh@Decisional Diffie-Hellman (DDH)|textbf}-Hellman (DDH) problem is hard for $\grgen$ if for any \ppt algorithm $\adv$,
  \[
  \advantage{DDH}{\adv, \grgen} \defeq \left|\pr{\game{\Game~\algo{DDH}}{\grgen} = 1} - \frac{1}{2}\right| = \negl.
  \]
\end{definition}

\subsection{ElGamal Commitments}

\begin{definition}[ElGamal Commitment Scheme]
  Let $\grgen$ be a group generation\index{group generation algorithm@Group generation algorithm} algorithm.
  The ElGamal\index{elgamal commitment@ElGamal commitment|textbf} commitment\index{commitment scheme@Commitment scheme} scheme $\algo{ElGamalCom}[\grgen]$ is defined as follows:
  \begin{itemize}
  \item $\algo{Setup}(1^\secpar) \rightarrow \params$: Run $\gparam \gets \grgen(1^\secpar)$, sample $h \sample \GG \setminus \{1_\GG\}$, and output $\params = (\gparam, h, \mathcal{M} \defeq \ZZ_p, \mathcal{R} \defeq \ZZ_p)$.
  \item $\algo{Commit}(\params, m, r) \rightarrow C$: Parse $\params = ((\GG, p, g), h, \mathcal{M}, \mathcal{R})$. For message $m \in \mathcal{M}$ and randomness $r \in \mathcal{R}$, output $C = (g^r, h^{m+r})$.
  \end{itemize}
\end{definition}

\begin{theorem}[ElGamal Commitments are Hiding]\label{thm:elgamal-hiding}
  Let $\grgen$ be a group generation algorithm for which the DDH\index{decisional diffiehellman ddh@Decisional Diffie-Hellman (DDH)} problem is hard, then the ElGamal commitment scheme $\algo{ElGamalCom}[\grgen]$ is hiding\index{hiding@Hiding}.
  More precisely, for any \ppt adversary $\adv$ against the hiding property of $\algo{ElGamalCom}[\grgen]$, there exists a \ppt adversary $\bdv$ against DDH such that
  \[
  \advantage{Hide}{\adv, \algo{ElGamalCom}} = 2 \cdot \advantage{DDH}{\bdv, \grgen}.
  \]
\end{theorem}

The proof is left as an exercise (see \autoref{ex:elgamal-hiding}).

\subsection{The OMDL and AOMDL Problems}

\begin{figure}[tbhp]
  \begin{center}
    \begin{tcolorbox}[width=\textwidth]
      \begin{pchstack}[center]
        \begin{pcvstack}
          \procedure[headlinesep=1pt]{$\game{\Game~\algo{OMDL}}{\grgen}$}{%
            \gparam \gets \grgen(1^\secpar) \\
            \ell \defeq 0 \pcsc q \defeq 0 \\
            \vect{y} \gets \adv^{\pcoracle{Chal}, \pcoracle{DL}}(\gparam) \\
            \vect{x} \defeq (x_1, \ldots, x_\ell) \\
            \pcreturn (\vect{y} = \vect{x} \wedge q < \ell)
          }
          \pcvspace
          \procedure[headlinesep=1pt]{Oracle $\pcoracle{DL}(X)$}{%
            q \defeq q+1 \\
            \pcreturn \log_g(X)
          }
        \end{pcvstack}
        \pchspace
        \begin{pcvstack}
          \procedure[headlinesep=1pt]{Oracle $\pcoracle{Chal}()$}{%
            \ell \defeq \ell + 1 \\
            x_\ell \sample \ZZ_p \\
            X_\ell \defeq g^{x_\ell} \\
            \pcreturn X_\ell
          }
        \end{pcvstack}
        \pchspace
        \begin{pcvstack}
          \procedure[headlinesep=1pt]{$\game{\Game~\algo{AOMDL}}{\grgen}$}{%
            \gparam \gets \grgen(1^\secpar) \\
            \ell \defeq 0 \pcsc q \defeq 0 \\
            \vect{y} \gets \adv^{\pcoracle{Chal}, \gamechange{$\oracle{ADL}$}}(\gparam) \\
            \vect{x} \defeq (x_1, \ldots, x_\ell) \\
            \pcreturn (\vect{y} = \vect{x} \wedge q < \ell)
          }
          \pcvspace
          \procedure[headlinesep=1pt]{Oracle $\oracle{ADL}((\alpha, \beta_1, \dots, \beta_{\ell}))$}{%
            q \defeq q+1 \\
            \pcreturn \alpha + \tsum_{i=1}^{\ell} {\beta_i x_i} \\
            \pclinecomment{${} = \log_g\left(g^\alpha \tprod_{i=1}^{\ell} {X_i^{\beta_i}}\right)$}
          }
        \end{pcvstack}
      \end{pchstack}
    \end{tcolorbox}
  \end{center}
  \caption{The one-more discrete logarithm (OMDL) and algebraic one-more discrete logarithm (AOMDL) games. The difference is \gamechange{highlighted}.\label{fig:omdl-aomdl-games}}
\end{figure}

\begin{definition}[One-More Discrete Logarithm Assumption]
  Let $\grgen$ be a group generation\index{group generation algorithm@Group generation algorithm} algorithm, and let $\game{\Game~\algo{OMDL}}{\grgen}$ be as defined in \autoref{fig:omdl-aomdl-games}.
  The one-more discrete logarithm\index{discrete logarithm dl@Discrete logarithm (DL)}\index{onemore discrete logarithm omdl@One-more discrete logarithm (OMDL)|textbf} (OMDL) problem is hard for $\grgen$ if for any \ppt algorithm $\adv$,
  \[
  \advantage{OMDL}{\adv, \grgen} \defeq \pr{\game{\Game~\algo{OMDL}}{\grgen} = 1} = \negl.
  \]
\end{definition}

\begin{definition}[Algebraic One-More Discrete Logarithm Assumption]
  Let $\grgen$ be a group generation algorithm, and let $\game{\Game~\algo{AOMDL}}{\grgen}$ be as defined in \autoref{fig:omdl-aomdl-games}.
  The algebraic one-more discrete logarithm\index{algebraic onemore discrete logarithm aomdl@Algebraic one-more discrete logarithm (AOMDL)|textbf} (AOMDL) problem is hard for $\grgen$ if for any \ppt algorithm $\adv$,
  \[
  \advantage{AOMDL}{\adv, \grgen} \defeq \pr{\game{\Game~\algo{AOMDL}}{\grgen} = 1} = \negl.
  \]
\end{definition}

\begin{remark}
  When proving the security of a scheme using the AOMDL assumption, the $\pcoracle{ADL}$ oracle is (typically) used by the reduction\index{reduction@Reduction} and not by the adversary against the scheme.
  Therefore, the security proof is not restricted to adversaries that know the algebraic representations of all group elements that they output.
\end{remark}

\begin{remark}[Falsifiable Assumptions]
  In order to be able to evaluate whether an assumption is true or not, it must be falsifiable\index{falsifiable assumption@Falsifiable assumption|textbf}, i.e., there must be a (constructive) way to demonstrate that it is false, if this is the case.
  More precisely, a cryptographic assumption is \emph{falsifiable} if there exists a \ppt challenger algorithm that interacts with an adversary and decides whether the adversary breaks the assumption.
\end{remark}

\subsection{Exercises}

\begin{exercise}
  In the framework of asymptotic security\index{asymptotic security@Asymptotic security}, why is the DL\index{discrete logarithm dl@Discrete logarithm (DL)} problem on secp256k1 not hard? What does $\grgen$ do?
\end{exercise}

\ifsolutions
\begin{mysolution}
  The DL\index{discrete logarithm dl@Discrete logarithm (DL)} problem on secp256k1 is not hard in the asymptotic security\index{asymptotic security@Asymptotic security} framework because secp256k1 has a fixed 256-bit group order, independent of the security parameter\index{security parameter@Security parameter} $\secpar$.
  For asymptotic hardness, the group size must grow with $\secpar$.
  The group generation\index{group generation algorithm@Group generation algorithm} algorithm $\grgen$ generates groups where $p \geq 2^{\secpar-1}$, ensuring that the problem scales appropriately with the security parameter.
\end{mysolution}
\fi

\begin{exercise}\label{ex:pedersen-binding}
  Argue that Pedersen\index{pedersen commitment@Pedersen commitment} commitments\index{commitment scheme@Commitment scheme} are binding\index{binding@Binding} under the DL\index{discrete logarithm dl@Discrete logarithm (DL)} assumption (\autoref{thm:pedersen-binding}).
\end{exercise}

\ifsolutions
\begin{mysolution}
  Let $\adv$ be an adversary that breaks the binding\index{binding@Binding} property of Pedersen\index{pedersen commitment@Pedersen commitment} commitments\index{commitment scheme@Commitment scheme}.
  That is, $\adv$ outputs $(m_0, r_0, m_1, r_1)$ such that $m_0 \neq m_1$ and $g^{m_0}h^{r_0} = g^{m_1}h^{r_1}$.
  We construct a \ppt reduction\index{reduction@Reduction} $\bdv$ that breaks DL\index{discrete logarithm dl@Discrete logarithm (DL)} as follows:
  \begin{itemize}
    \item $\bdv$ receives $(\gparam, h)$ where $h = g^z$ for unknown $z$.
    \item $\bdv$ runs $\adv$ with parameters $\params = (\gparam, h)$.
    \item When $\adv$ outputs $(m_0, r_0, m_1, r_1)$, we have $g^{m_0}h^{r_0} = g^{m_1}h^{r_1}$.
    \item This implies $g^{m_0 - m_1} = h^{r_1 - r_0}$, so $g^{m_0 - m_1} = g^{z(r_1 - r_0)}$.
    \item If $r_0 = r_1$, then $g^{m_0 - m_1} = 1$, which contradicts $m_0 \neq m_1$ (since $g$ generates $\GG$).
    \item Therefore $r_0 \neq r_1$, and we can compute $z = \frac{m_0 - m_1}{r_1 - r_0} \bmod p$.
    \item $\bdv$ outputs $z$.
  \end{itemize}
  The reduction is perfect: whenever $\adv$ succeeds in breaking binding, $\bdv$ succeeds in computing the discrete log.
  Therefore, $\advantage{DL}{\bdv, \grgen} = \advantage{Bind}{\adv, \algo{PedCom}}$.
\end{mysolution}
\fi

\begin{exercise}[Optional]
  Knowledge of the discrete logarithm\index{discrete logarithm dl@Discrete logarithm (DL)} $x$ such that $g^x = h$ allows efficiently producing a Pedersen\index{pedersen commitment@Pedersen commitment} commitment\index{commitment scheme@Commitment scheme} $C$ such that $\algo{Commit}(\params, m, r) = \algo{Commit}(\params, mx \bmod p, r - 1 \bmod p)$.
  Why does this not violate \autoref{thm:pedersen-binding}?
  
  Hint: We have the relationship $g^m h^r = g^m (g^x)^r = g^{m + xr} = g^{mx} (g^x)^{r-1} = g^{mx} h^{r-1}$.
\end{exercise}

\ifsolutions
\begin{mysolution}
  This does not violate the binding\index{binding@Binding} theorem because knowing $x = \log_g h$ means the DL\index{discrete logarithm dl@Discrete logarithm (DL)} problem has been solved for this instance.
  The binding theorem only claims that \emph{if} the DL problem is hard, \emph{then} the commitment\index{commitment scheme@Commitment scheme} scheme is binding.
  Since the hypothesis is not satisfied, the theorem makes no claim about binding.
  
  This highlights that Pedersen\index{pedersen commitment@Pedersen commitment} commitments rely crucially on the setup assumption that no one knows the discrete logarithm relationship between $g$ and $h$.
\end{mysolution}
\fi

\begin{exercise}\label{ex:pedersen-hiding}
  Argue that Pedersen\index{pedersen commitment@Pedersen commitment} commitments\index{commitment scheme@Commitment scheme} are perfectly hiding\index{hiding@Hiding}\index{hiding@Hiding!perfectly hiding@perfectly hiding} (\autoref{thm:pedersen-hiding}).
\end{exercise}

\ifsolutions
\begin{mysolution}
  For any message $m \in \ZZ_p$, the commitment\index{commitment scheme@Commitment scheme} $C = g^m h^r$ is uniformly distributed over $\GG$ when $r$ is chosen uniformly from $\ZZ_p$.
  Since $h \in \GG \setminus \{1_\GG\}$ and $\GG$ has prime order $p$, $h$ generates $\GG$ (every non-identity element generates a prime-order group).
  Thus the map $r \mapsto h^r$ is a bijection from $\ZZ_p$ to $\GG$.
  For fixed $m$, the map $r \mapsto g^m h^r$ is also a bijection from $\ZZ_p$ to $\GG$.
  Therefore, for any two messages $m_0, m_1$, the distributions of $\algo{Commit}(\params, m_0, r)$ and $\algo{Commit}(\params, m_1, r)$ for uniform $r$ are both uniform over $\GG$.
\end{mysolution}
\fi

\begin{exercise}[Optional]\label{ex:elgamal-hiding}
  Prove that ElGamal\index{elgamal commitment@ElGamal commitment} commitments\index{commitment scheme@Commitment scheme} are hiding\index{hiding@Hiding} under the DDH\index{decisional diffiehellman ddh@Decisional Diffie-Hellman (DDH)} assumption (\autoref{thm:elgamal-hiding}).
\end{exercise}

\ifsolutions
\begin{mysolution}
   Let $\adv$ be an adversary against the hiding\index{hiding@Hiding} property of ElGamal\index{elgamal commitment@ElGamal commitment}.
  We construct a \ppt reduction\index{reduction@Reduction} $\bdv$ against DDH\index{decisional diffiehellman ddh@Decisional Diffie-Hellman (DDH)}:
  \begin{itemize}
    \item $\bdv$ receives $(\gparam, A, B, C)$ where $A = g^a$, $B = g^b$, and either $C = g^{ab}$ or $C = g^c$ for random $c$.
    \item $\bdv$ sets $h = B$ and gives $\params = (\gparam, h)$ to $\adv$.
    \item When $\adv$ outputs $(m_0, m_1)$, $\bdv$ chooses $e \sample \{0, 1\}$ and computes the commitment\index{commitment scheme@Commitment scheme} $(A, B^{m_e} \cdot C)$.
    \item $\bdv$ gives this commitment to $\adv$.
    \item When $\adv$ outputs $e'$, $\bdv$ outputs $1$ if $e' = e$ and $0$ otherwise.
  \end{itemize}

  Analysis:
  \begin{itemize}
      \item When $C = g^{ab}$ (real tuple): The commitment $(A, B^{m_e} \cdot C) = (g^a, g^{b \cdot m_e} \cdot g^{ab}) = (g^a, g^{b(m_e + a)}) = (g^a, h^{m_e + a})$.
      This is a valid ElGamal commitment to $m_e$ with randomness $r = a$.
      Thus $\bdv$ perfectly simulates the hiding game, so $\left|\Pr[e = e' | C = g^{ab}] - \frac{1}{2}\right| = \advantage{Hide}{\adv, \algo{ElGamalCom}}$.
      \item When $C = g^c$ (random tuple): The commitment is $(A, B^{m_e} \cdot C) = (g^a, g^{bm_e + c})$.
      Since $c$ is uniformly random and independent of everything else, the second component is uniformly random in $\GG$.
      The commitment reveals no information about $e$, so $\Pr[e = e' | C = g^c] = \frac{1}{2}$.
  \end{itemize}
  For the DDH advantage\index{advantage@Advantage}, let $d \in \{0, 1\}$ be the bit indicating whether the DDH tuple is real ($d = 1$ if $C = g^{ab}$, $d = 0$ if $C = g^c$), and let $d'$ be $\bdv$'s output bit.
  Recall that $\bdv$ outputs $d' = 1$ if $e' = e$ and $d' = 0$ otherwise.
  \begin{align*}
  \advantage{DDH}{\bdv, \grgen} &= \left|\Pr[d' = d] - \frac{1}{2}\right|\\
  \Pr[d' = d] &= \Pr[d' = 1 \wedge d = 1] + \Pr[d' = 0 \wedge d = 0]\\
  &= \Pr[d' = 1 | d = 1] \cdot \Pr[d = 1] + \Pr[d' = 0 | d = 0] \cdot \Pr[d = 0]\\
  &= \frac{1}{2} \Pr[d' = 1 | C = g^{ab}] + \frac{1}{2} \Pr[d' = 0 | C = g^c]\\
  &= \frac{1}{2} \Pr[d' = 1 | C = g^{ab}] + \frac{1}{2}(1 - \Pr[d' = 1 | C = g^c])\\
  &= \frac{1}{2} \Pr[d' = 1 | C = g^{ab}] + \frac{1}{2} - \frac{1}{2}\Pr[d' = 1 | C = g^c]
  \end{align*}
  Therefore:
  \begin{align*}
  \advantage{DDH}{\bdv, \grgen} &= \left|\frac{1}{2} \Pr[d' = 1 | C = g^{ab}] - \frac{1}{2}\Pr[d' = 1 | C = g^c]\right|\\
  &= \frac{1}{2} \left|\Pr[d' = 1 | C = g^{ab}] - \Pr[d' = 1 | C = g^c]\right|\\
  &= \frac{1}{2} \left|\Pr[e = e' | C = g^{ab}] - \Pr[e = e' | C = g^c]\right|\\
  &= \frac{1}{2} \left|\Pr[e = e' | C = g^{ab}] - \frac{1}{2}\right|\\
  &= \frac{1}{2} \cdot \advantage{Hide}{\adv, \algo{ElGamalCom}}
  \end{align*}

  This gives us: $\advantage{Hide}{\adv, \algo{ElGamalCom}} = 2 \cdot \advantage{DDH}{\bdv, \grgen}$. $\bdv$ is \ppt\ if $\adv$ is.
\end{mysolution}
\fi

\begin{exercise}
  Argue that if OMDL\index{onemore discrete logarithm omdl@One-more discrete logarithm (OMDL)} is hard for a group generation\index{group generation algorithm@Group generation algorithm} algorithm then DL\index{discrete logarithm dl@Discrete logarithm (DL)} is hard.
\end{exercise}

\ifsolutions
\begin{mysolution}
  If OMDL\index{onemore discrete logarithm omdl@One-more discrete logarithm (OMDL)} is hard, then DL\index{discrete logarithm dl@Discrete logarithm (DL)} must be hard.
  Given a \ppt DL adversary $\adv$, we construct a \ppt OMDL adversary $\bdv$:
  \begin{itemize}
    \item $\bdv$ calls $\pcoracle{Chal}()$ once to get $X_1$.
    \item $\bdv$ runs $\adv(\gparam, X_1)$ to get $x_1'$.
    \item $\bdv$ outputs $(x_1')$.
  \end{itemize}
  If $\adv$ succeeds, then $X_1 = g^{x_1'}$, so $\bdv$ solves one discrete log with zero challenge queries, winning the OMDL game.
\end{mysolution}
\fi

\begin{exercise}
  Consider an algorithm $\adv$ playing the AOMDL\index{algebraic onemore discrete logarithm aomdl@Algebraic one-more discrete logarithm (AOMDL)} game. It queries $\pcoracle{Chal}$ twice to obtain $X_1$ and $X_2$, then samples $\alpha, \beta_1, \beta_2 \sample \ZZ_p$. 
  \begin{enumerate}
    \item How should $\adv$ call the $\oracle{ADL}$ oracle to obtain $\log_g(P)$ where $P = g^\alpha X_1^{\beta_1} X_2^{\beta_2}$?
    \item Assuming $\adv$ makes no more $\pcoracle{Chal}$ queries, how many more $\oracle{ADL}$ queries can it make?
    \item What must $\adv$ return to win the game?
  \end{enumerate}
\end{exercise}

\ifsolutions
\begin{mysolution}
  \begin{enumerate}
  \item $\adv$ should call $\oracle{ADL}((\alpha, \beta_1, \beta_2))$. The oracle will return
    \begin{align*}
      \alpha + \beta_1 x_1 + \beta_2 x_2 &= \log_g(g^\alpha X_1^{\beta_1} X_2^{\beta_2})\\
      &= \log_g(P)
    \end{align*}.
    \item Since $\ell = 2$ (two $\pcoracle{Chal}$ queries) and $q = 1$ (one $\oracle{ADL}$ query), and the winning condition requires $q < \ell$, the adversary can make at most 0 more queries (as $1 < 2$ but $2 \not< 2$).
    \item $\adv$ must return $(x_1, x_2)$, the discrete logarithms\index{discrete logarithm dl@Discrete logarithm (DL)} of $X_1$ and $X_2$.
  \end{enumerate}
\end{mysolution}
\fi

\begin{exercise}
  Is the OMDL\index{onemore discrete logarithm omdl@One-more discrete logarithm (OMDL)} assumption falsifiable\index{falsifiable assumption@Falsifiable assumption}? Is the AOMDL\index{algebraic onemore discrete logarithm aomdl@Algebraic one-more discrete logarithm (AOMDL)} assumption falsifiable? Do we prefer the security of a scheme to be based on OMDL or AOMDL?
\end{exercise}

\ifsolutions
\begin{mysolution}
  The OMDL\index{onemore discrete logarithm omdl@One-more discrete logarithm (OMDL)} assumption is \emph{non-falsifiable\index{falsifiable assumption@Falsifiable assumption}} because implementing the DL\index{discrete logarithm dl@Discrete logarithm (DL)} oracle requires computing discrete logarithms, but the OMDL assumption itself assumes that this cannot be done with non-negligible\index{negligible function@Negligible function} success probability in polynomial time\index{algorithm@Algorithm!polynomialtime@polynomial-time} classically.

  The AOMDL\index{algebraic onemore discrete logarithm aomdl@Algebraic one-more discrete logarithm (AOMDL)} assumption is \emph{falsifiable} because:
  \begin{itemize}
    \item The ADL oracle can be implemented efficiently: given $(\alpha, \beta_1, \ldots, \beta_\ell)$, it simply returns $\alpha + \sum_{i=1}^\ell \beta_i x_i$
    \item The winning condition can be verified by checking that the returned values equal the stored $x_i$ values
    \item All operations (additions, multiplications) are polynomial time
  \end{itemize}

  We prefer basing security on falsifiable assumptions like AOMDL. Note that AOMDL is a weaker assumption than OMDL---if AOMDL doesn't hold, then OMDL doesn't hold either (since any OMDL adversary is also an AOMDL adversary).
\end{mysolution}
\fi

\begin{exercise}[Optional]
  Why is the DL\index{discrete logarithm dl@Discrete logarithm (DL)} problem not hard in general?
\end{exercise}

\ifsolutions
\begin{mysolution}
  The discrete logarithm\index{discrete logarithm dl@Discrete logarithm (DL)} in $\ZZ_p$ where the group operation is addition can be computed in polynomial time\index{algorithm@Algorithm!polynomialtime@polynomial-time}.
  In general, quantum\index{postquantum security@Post-quantum security} computers can solve the DL problem in polynomial time using Shor's algorithm.
\end{mysolution}
\fi

\begin{exercise}[Optional]
  Read about the Generic Group Model\index{generic group model ggm@Generic group model (GGM)|textbf} (GGM) in the introduction of~\cite{EC:Shoup97} and~\cite{IMA:Maurer05}, and argue why the DLP\index{discrete logarithm dl@Discrete logarithm (DL)} is hard for classical computational models.
\end{exercise}

\begin{exercise}[Optional]
  How does Silent Payments\index{silent payments@Silent Payments}~\cite{add:bip-silentpayments} rely on the DDH\index{decisional diffiehellman ddh@Decisional Diffie-Hellman (DDH)} assumption?
\end{exercise}

\begin{exercise}[Optional]
  Study the concept of falsifiability\index{falsifiable assumption@Falsifiable assumption} of cryptographic assumptions by reading Naor~\cite{C:Naor03}.
  Analyze the following assumptions and explain the extent to which each one is falsifiable:
  \begin{itemize}
    \item Factoring
    \item The RSA assumption
    \item The Decisional Diffie-Hellman assumption
    \item The Knowledge-of-Exponent assumption
  \end{itemize}
\end{exercise}
